\Appendix{Kode Sumber Sistem Retrieval-Augmented Generation (RAG)}
\label{lampiran:kode_rag}

Pada lampiran ini, dilampirkan kode sumber dari sistem \textit{Retrieval-Augmented Generation} (RAG) \textit{version-aware} yang telah dikembangkan. Kode sumber dibagi menjadi beberapa modul utama berdasarkan fungsinya, yaitu pemrosesan dokumen, layanan \textit{backend} API, antarmuka pengguna berbasis teks, serta skrip evaluasi sistem.

\section{Pemrosesan Dokumen dan Basis Data Vektor}

Bagian ini memuat skrip yang bertanggung jawab untuk mempersiapkan pangkalan pengetahuan (\textit{knowledge base}). Skrip \texttt{ingest.py} berfungsi untuk mengekstraksi teks dari berbagai format dokumen (PDF, gambar, maupun teks biasa), membaginya menjadi potongan-potongan informasi (\textit{chunks}) dengan tetap mempertahankan konteks lintas halaman, lalu mengubahnya menjadi representasi vektor (\textit{embeddings}) menggunakan model \texttt{intfloat/multilingual-e5-base} sebelum disimpan ke dalam ChromaDB.

\begin{myblock}{Skrip \textit{Ingestion} Dokumen (\texttt{ingest.py})}\label{kode:ingest}
\begin{minted}[xleftmargin=20pt,linenos,breaklines]{python}

def ingest_file(path, version, collection, model, manual_name=None):
    WINDOW_SIZE = 15   
    STRIDE = 5        
    
    i = 0
    while i < len(global_lines):
        window_items = global_lines[i : i + WINDOW_SIZE]
        chunk_text = "\n".join([item['text'] for item in window_items]).strip()
        start_page = window_items[0]['page']
        end_page = window_items[-1]['page']
        page_str = str(start_page) if start_page == end_page else f"{start_page}-{end_page}"
        
        if len(chunk_text) > 30:
            meta = {
                "doc_id": doc_uuid,
                "name": name,             
                "version": version,
                "timestamp": ts,
                "page": page_str,
                "content": chunk_text
            }
            all_chunk_docs.append(chunk_text)
            all_chunk_metas.append(meta)
            all_chunk_ids.append(
                f"{doc_uuid}_{total_chunks}")
            total_chunks += 1
        i += STRIDE

    embeddings = model.encode(all_chunk_docs, normalize_embeddings=True)
    collection.add(
        documents=all_chunk_docs,
        embeddings=embeddings.tolist(),
        metadatas=all_chunk_metas,
        ids=all_chunk_ids
    )
\end{minted}
\end{myblock}

Selanjutnya, skrip \texttt{inspect\_db.py} digunakan sebagai alat bantu pemantauan untuk memverifikasi isi dari ChromaDB. Skrip ini membaca metadata dari pangkalan data dan menampilkan informasi ringkas mengenai daftar dokumen, versi dokumen, jumlah \textit{chunk}, dan stempel waktu (\textit{timestamp}).

\begin{myblock}{Skrip Inspeksi Basis Data (\texttt{inspect\_db.py})}\label{kode:inspect}
\begin{minted}[xleftmargin=20pt,linenos,breaklines]{python}

def main():
    client = chromadb.PersistentClient(path="./chroma_db")
    collection = client.get_collection("docs_versioned")
    all_data = collection.get(include=["metadatas"])
    
    # Struktur data untuk mengelompokkan chunk berdasarkan nama file
    files_summary = defaultdict(lambda: {
        'version': 'Unknown', 'timestamp': 'N/A', 
        'doc_id': 'N/A', 'chunks': 0, 'pages': set()
    })

    for meta in all_data['metadatas']:
        fname = meta.get('filename') or meta.get('source') or "Tanpa Nama"
        clean_fname = fname.split('/')[-1].split('\\')[-1]

        files_summary[clean_fname]['chunks'] += 1
        
        # Ambil metadata umum pada temuan potongan teks (chunk) pertama
        if files_summary[clean_fname]['chunks'] == 1:
            files_summary[clean_fname]['version'] = meta.get('version', 'N/A')
            files_summary[clean_fname]['timestamp'] = meta.get('timestamp', 'N/A')
        
        # Hitung set halaman unik untuk memastikan kelengkapan parsing dokumen
        if 'page' in meta and meta['page'] is not None:
            files_summary
            [clean_fname]['pages'].add(meta['page'])

\end{minted}
\end{myblock}

\section{Layanan \textit{Backend} API RAG}

Skrip \texttt{app.py} merupakan inti dari sistem RAG yang dibangun menggunakan \textit{framework} FastAPI. Modul ini bertugas menerima permintaan kueri dari pengguna, melakukan pencarian semantik (\textit{similarity search}) pada ChromaDB dengan mempertimbangkan spesifikasi versi dokumen (\textit{version-aware}), merakit \textit{prompt} yang komprehensif, dan mengirimkannya ke model bahasa (LLM) Ollama untuk menghasilkan respons akhir berbahasa Indonesia.

\begin{myblock}{Skrip \textit{Backend} API (\texttt{app.py})}\label{kode:app}
\begin{minted}[xleftmargin=20pt,linenos,breaklines]{python}

@app.post("/query", response_model=QueryResponse)
async def handle_query(request: QueryRequest):
    if request.version == "latest":
        # Ambil sampel lebih besar untuk memfilter dokumen terbaru
        results = collection.query(
            query_embeddings=[q_emb_list], n_results=TOP_K_SEARCH,
            include=["metadatas", "distances", "documents"]
        )
        
        # Pengelompokan dokumen berdasarkan nama
        doc_groups = {}
        for (doc, meta, dist) in zip(results['documents'][0], results['metadatas'][0], results['distances'][0]):
            group_key = meta.get("name")
            if group_key not in doc_groups: doc_groups[group_key] = []
            doc_groups[group_key].append({"chunk": doc, "meta": meta, "dist": dist})

        final_candidates = []
        for name, items in doc_groups.items():
            # Filter: Hanya ambil potongan dokumen dengan timestamp terbaru
            latest_ts = max(item['meta'].get('timestamp', 0) for item in items)
            latest_chunks = [item for item in items if item['meta'].get('timestamp', 0) == latest_ts]
            final_candidates.extend(latest_chunks)
            
        final_candidates.sort(key=lambda x: x["dist"]) 
        top_chunks_with_meta = final_candidates[:request.top_k]

    else:
        results = collection.query(
            query_embeddings=[q_emb_list], n_results=request.top_k,
            where={"version": request.version},
            include=["metadatas", "distances", "documents"]
        )

    prompt = build_prompt(request.query, top_chunks_with_meta)
    answer = call_ollama_chat(prompt)
    
    return QueryResponse(answer=answer, provenance=format_provenance(top_chunks_with_meta))
\end{minted}
\end{myblock}

\section{Antarmuka Pengguna Berbasis Baris Perintah}

Untuk memfasilitasi interaksi pengguna dengan layanan \textit{backend}, dikembangkan sebuah antarmuka berbasis baris perintah (\textit{Command Line Interface}) pada skrip \texttt{rag\_cli\_versionaware.py}. Antarmuka ini memungkinkan pengguna untuk mengirimkan kueri, mengganti versi dokumen target yang ingin dicari, serta menampilkan jawaban LLM beserta tabel sumber dokumen asal (\textit{provenance}) dan skor kesamaannya.

\begin{myblock}{Skrip Antarmuka CLI (\texttt{rag\_cli\_versionaware.py})}\label{kode:cli}
\begin{minted}[xleftmargin=20pt,linenos,breaklines]{python}

def ask_rag(query: str, version: str = "latest"):
    # Mengirimkan kueri beserta parameter versi ke endpoint RAG backend
    payload = {"query": query, "version": version, "top_k": 5}
    res = requests.post("http://localhost:8000/query", json=payload, timeout=120)
    return res.json() if res.status_code == 200 else None

def main():
    version = "latest" # Konfigurasi versi pencarian bawaan
    
    while True:
        cmd = Prompt.ask("[bold green] Pertanyaan / Perintah[/bold green]").strip()
        if cmd.lower() in ["exit", "quit"]: break

        if cmd.lower().startswith("pakai "):
            version = cmd.split("pakai ", 1)[1].strip()
            console.print(f" Versi diubah menjadi: {version}")
            continue

        result = ask_rag(cmd, version)
        if not result: continue

        answer = result.get("answer", "").strip()
        provenance = result.get("provenance", [])

        # Menampilkan jawaban utama dan tabel sumber rujukan
        console.print(Panel.fit(answer, title=f" Jawaban ({version})"))
        if provenance:
            show_provenance_table(provenance)
        
        log_chat(cmd, version, answer, provenance)
\end{minted}
\end{myblock}

\section{Evaluasi Kinerja Sistem}

Tahap pengujian kinerja sistem RAG diimplementasikan dalam skrip \texttt{eval\_ragas.py}. Skrip ini menggunakan kerangka kerja RAGAS (\textit{Retrieval Augmented Generation Assessment}) untuk menghitung metrik evaluasi seperti \textit{faithfulness}, \textit{answer relevancy}, \textit{context precision}, dan \textit{context recall}. Hasil evaluasi ini kemudian diekspor ke dalam format \textit{spreadsheet} untuk analisis lebih lanjut pada Bab 4.

\begin{myblock}{Skrip Evaluasi Kinerja Sistem (\texttt{eval\_ragas.py})}\label{kode:eval}
\begin{minted}[xleftmargin=20pt,linenos,breaklines]{python}

def run_evaluation():
    judge_llm = ChatOllama(model="llama3:8b", temperature=0, format="json")
    judge_embeddings = HuggingFaceEmbeddings(
        model_name="intfloat/multilingual-e5-base") 
    
    # Konversi data inferensi menjadi Dataset HuggingFace
    dataset = Dataset.from_dict(data_dict)
    results = evaluate(
        dataset=dataset,
        metrics=[
            faithfulness,
            answer_relevancy,
            context_precision,
            context_recall,
        ],
        llm=judge_llm,
        embeddings=judge_embeddings,
        run_config=RunConfig(max_workers=1, timeout=300)
    )

    # Simpan laporan evaluasi ke spreadsheet
    df = results.to_pandas()
    df.to_excel("laporan_evaluasi_ragas.xlsx", index=False)
\end{minted}
\end{myblock}